% Libro 2 · Mundo 2 · Lógica proposicional (app: pu-logica)
\mundo{Lógica proposicional}{pu-logica}{La lógica es la gramática del razonamiento: dice cuándo una conclusión se
sigue de unas premisas y cuándo solo lo parece. En el examen aparece de forma explícita (tablas de verdad, negaciones,
silogismos) y de forma escondida en casi todas las demás preguntas, cada vez que hay un «si…, entonces…» o un
«todos».}{%
\item \textbf{2.1} Conectivos y tablas de verdad
\item \textbf{2.2} Negación y leyes de De Morgan
\item \textbf{2.3} El condicional
\item \textbf{2.4} Cuantificadores y silogismos
\item \textbf{2.5} Verdad y mentira}

% ------------------------------------------------------------------
\leccion{pu-lo-1}{Conectivos y tablas de verdad}{Juego de los interruptores que encienden una lámpara.}

\begin{reto}
Si $p$ es verdadera y $q$ es falsa, ¿cuál de las siguientes proposiciones es verdadera?
\opciones{$p \wedge q$}{$\neg p \vee q$}{$p \to q$}{$q \to p$}
\end{reto}

\subsection{La idea: circuitos que encienden}
Una \clave{proposición} es una frase que es verdadera (V) o falsa (F): «7 es primo» lo es; «¿qué hora es?» no. Los
\clave{conectivos} combinan proposiciones, y funcionan como interruptores:

\begin{itemize}
  \item $p \wedge q$ («$p$ y $q$»): interruptores en serie. Solo enciende si los dos están cerrados.
  \item $p \vee q$ («$p$ o $q$»): en paralelo. Enciende si al menos uno está cerrado.
  \item $\neg p$ («no $p$»): invierte.
  \item $p \to q$ («si $p$, entonces $q$»): una promesa. Solo se rompe si se cumple $p$ y no $q$.
  \item $p \leftrightarrow q$ («$p$ si y solo si $q$»): verdadera cuando las dos tienen el mismo valor.
\end{itemize}

\begin{center}\small
\begin{tabular}{cc|ccccc}
$p$ & $q$ & $p \wedge q$ & $p \vee q$ & $p \to q$ & $p \leftrightarrow q$ & $p \veebar q$\\\hline
V & V & V & V & V & V & F\\
V & F & F & V & \textbf{F} & F & V\\
F & V & F & V & V & F & V\\
F & F & F & F & V & V & F\\
\end{tabular}
\end{center}

\subsubsection{¿Por qué $F \to$ cualquier cosa es V?}
Piensa en la promesa «si ganas el examen, te invito a cine». Si no ganas, la promesa no se rompe, haya cine o no. El
condicional solo es falso cuando el antecedente se cumple y el consecuente no.

\subsubsection{Tautologías y contradicciones}
Una \clave{tautología} es verdadera en todas las filas ($p \vee \neg p$); una \clave{contradicción}, falsa en todas
($p \wedge \neg p$). Con $n$ proposiciones simples, la tabla tiene $2^n$ filas.

\figura{circuitos}{La conjunción y la disyunción como circuitos.}

\begin{regla}
\begin{itemize}
  \item $\wedge$: V solo si ambas V. $\vee$: F solo si ambas F.
  \item $\to$: F solo en el caso V $\to$ F.
  \item $\leftrightarrow$: V si coinciden. $\veebar$ (o exclusivo): V si difieren.
  \item Una tabla con $n$ variables tiene $2^n$ filas.
\end{itemize}
\end{regla}

\subsection{Ejemplos resueltos}
\begin{ejemplo}[Evaluar paso a paso]
Si $p$ es V, $q$ es F y $r$ es V, ¿cuál es el valor de $(p \vee q) \to (q \wedge r)$?
\solucion $p \vee q = $ V. $q \wedge r = $ F. V $\to$ F $=$ F.
\end{ejemplo}

\begin{ejemplo}[Deducir valores]
Se sabe que $(p \wedge q) \to r$ es falsa. ¿Qué valores tienen $p$, $q$ y $r$?
\solucion Un condicional falso exige antecedente V y consecuente F: $p \wedge q$ es V (así que $p$ y $q$ son V) y $r$
es F.
\end{ejemplo}

\begin{ejemplo}[¿Tautología?]
¿Es $(p \to q) \vee (q \to p)$ una tautología?
\solucion Para que sea falsa, ambos condicionales serían falsos: $p$ V y $q$ F en el primero, $q$ V y $p$ F en el
segundo. Imposible a la vez: es una tautología.
\end{ejemplo}

\begin{trampa}
Leer el «o» del español como exclusivo. En lógica, $p \vee q$ es verdadera también cuando las dos lo son. El «o» exclusivo
se escribe $\veebar$ y el examen lo dice expresamente («o…, pero no ambos»).
\end{trampa}

\begin{atajo}
Para saber cuándo una proposición es falsa, busca la única fila que la hace falsa (en $\to$ y $\vee$) en vez de hacer la
tabla completa.
\end{atajo}

\taller
\piso{Calienta}{evaluar}
\begin{ejercicios}[2]
\item $p$ V, $q$ V: $\neg p \to q$. \sol{F $\to$ V $=$ V.}
\item $p$ F, $q$ V: $p \leftrightarrow q$. \sol{F.}
\item $p$ V, $q$ F: $\neg(p \wedge \neg q)$. \sol{$p \wedge \neg q = $ V; negada: F.}
\item Filas de la tabla con 4 variables. \sol{16.}
\end{ejercicios}
\piso{Entrena}{deducir}
\begin{ejercicios}
\item Si $p \vee q$ es F, ¿qué valor tiene $p \to q$? \sol{$p$ y $q$ son F; F $\to$ F $=$ V.}
\item Si $\neg(p \to q)$ es V, ¿qué valor tiene $q \to p$? \sol{$p$ V, $q$ F: F $\to$ V $=$ V.}
\item ¿Es $p \wedge \neg p$ tautología, contradicción o contingencia? \sol{Contradicción.}
\end{ejercicios}
\piso{Como en el examen}{opción múltiple}
\begin{ejercicios}
\item Si $p \to q$ es falsa, entonces es verdadera
\opciones{$p \wedge q$}{$\neg p$}{$p \vee q$}{$q$}
\sol{C. $p$ es V y $q$ es F; $p \vee q$ es V.}
\item ¿Cuál es una tautología?
\opcionesd{$p \to (p \vee q)$}{$(p \vee q) \to p$}{$p \wedge \neg q$}{$p \leftrightarrow \neg p$}
\sol{A. Si $p$ es V, $p \vee q$ también.}
\end{ejercicios}

\begin{retoresuelto}
$p \wedge q$: F. $\neg p \vee q$: F $\vee$ F $=$ F. $p \to q$: V $\to$ F $=$ F. $q \to p$: F $\to$ V $=$ V.
\textbf{Opción D.}
\end{retoresuelto}
\bitacora

% ------------------------------------------------------------------
\leccion{pu-lo-2}{Negación y De Morgan}{Juego del espejo que niega frases.}

\begin{reto}
La negación de «Todos los estudiantes aprobaron y nadie se quejó» es
\opcionesl{Algún estudiante no aprobó o alguien se quejó.}{Ningún estudiante aprobó y todos se quejaron.}{Algún estudiante no aprobó y alguien se quejó.}{Todos los estudiantes reprobaron o alguien se quejó.}
\end{reto}

\subsection{La idea: negar es decir exactamente lo contrario}
La negación de una frase es la que es verdadera justo cuando la original es falsa. No es «lo opuesto extremo»: la negación
de «hace calor» es «no hace calor», no «hace frío».

\subsubsection{Leyes de De Morgan}
Para negar un «y», basta que falle una de las partes; para negar un «o», tienen que fallar las dos:
\[\neg(p \wedge q) \equiv \neg p \vee \neg q \qquad\qquad \neg(p \vee q) \equiv \neg p \wedge \neg q\]
«No es cierto que estudie y trabaje» equivale a «no estudia o no trabaja».

\subsubsection{Negar cuantificadores}
\begin{center}\small
\begin{tabular}{@{}ll@{}}
\toprule
\textbf{Frase} & \textbf{Negación}\\\midrule
Todos son $P$ & Algunos no son $P$ (al menos uno no es $P$)\\
Algunos son $P$ & Ninguno es $P$\\
Ninguno es $P$ & Algunos son $P$\\
Algunos no son $P$ & Todos son $P$\\\bottomrule
\end{tabular}
\end{center}

\subsubsection{Negar un condicional}
$\neg(p \to q) \equiv p \wedge \neg q$. La negación de «si llueve, llevo paraguas» es «llueve y no llevo paraguas»
(no «si no llueve, no llevo paraguas»).

\begin{regla}
\begin{itemize}
  \item $\neg(p \wedge q) \equiv \neg p \vee \neg q$; $\neg(p \vee q) \equiv \neg p \wedge \neg q$.
  \item $\neg(\text{todos } P) \equiv$ algunos no $P$; $\neg(\text{algunos } P) \equiv$ ninguno $P$.
  \item $\neg(p \to q) \equiv p \wedge \neg q$. $\neg\neg p \equiv p$.
\end{itemize}
\end{regla}

\subsection{Ejemplos resueltos}
\begin{ejemplo}[Una negación doble]
Niega: «Algunos médicos no usan gafas».
\solucion La negación de «algunos no son» es «todos son»: «Todos los médicos usan gafas».
\end{ejemplo}

\begin{ejemplo}[De Morgan en contexto]
Niega: «Juan habla inglés o francés».
\solucion «Juan no habla inglés y no habla francés» (no habla ninguno).
\end{ejemplo}

\begin{ejemplo}[Condicional]
Niega: «Si un número es par, entonces es divisible entre 4».
\solucion «Hay un número par que no es divisible entre 4» (por ejemplo, 6). En lógica proposicional: $p \wedge \neg q$.
\end{ejemplo}

\begin{trampa}
Negar «todos aprobaron» con «ninguno aprobó». Basta que uno no haya aprobado para que «todos aprobaron» sea falsa.
\end{trampa}

\taller
\piso{Entrena}{negar}
\begin{ejercicios}
\item Niega: «Ningún perro maúlla». \sol{«Algún perro maúlla».}
\item Niega: «Estudio y no veo televisión». \sol{«No estudio o veo televisión».}
\item Niega: «Si como, no duermo». \sol{«Como y duermo».}
\item ¿Es «no todos vinieron» lo mismo que «algunos no vinieron»? \sol{Sí.}
\end{ejercicios}
\piso{Como en el examen}{opción múltiple}
\begin{ejercicios}
\item La negación de «Algunos números primos son pares» es
\opcionesl{Algunos números primos no son pares.}{Ningún número primo es par.}{Todos los números primos son pares.}{Todos los números pares son primos.}
\sol{B. (Y como 2 es primo y par, la original es verdadera y la negación falsa.)}
\item $\neg(\neg p \vee q)$ equivale a
\opciones{$p \wedge \neg q$}{$\neg p \wedge q$}{$p \vee \neg q$}{$p \to q$}
\sol{A. De Morgan: $p \wedge \neg q$.}
\end{ejercicios}

\begin{retoresuelto}
La frase es «$A$ y $B$» con $A$ = «todos aprobaron» y $B$ = «nadie se quejó». Su negación es «$\neg A$ o $\neg B$»:
«algún estudiante no aprobó o alguien se quejó». \textbf{Opción A.}
\end{retoresuelto}
\bitacora

% ------------------------------------------------------------------
\leccion{pu-lo-3}{El condicional}{Juego de las cartas que hay que voltear.}

\begin{reto}
«Si llueve, el partido se cancela». Se sabe que el partido no se canceló. ¿Qué se puede concluir?
\opcionesd{Llovió.}{No llovió.}{No se sabe si llovió.}{Se jugó bajo la lluvia.}
\end{reto}

\subsection{La idea: cuatro frases que se parecen}
A partir de «si $p$, entonces $q$» se forman tres frases parecidas. Solo una dice lo mismo:

\begin{center}\small
\begin{tabular}{@{}lll@{}}
\toprule
\textbf{Nombre} & \textbf{Forma} & \textbf{¿Equivale a $p \to q$?}\\\midrule
Condicional & si $p$, entonces $q$ & ---\\
Recíproco & si $q$, entonces $p$ & no\\
Inverso & si no $p$, entonces no $q$ & no\\
Contrarrecíproco & si no $q$, entonces no $p$ & \textbf{sí}\\\bottomrule
\end{tabular}
\end{center}

«Si es un perro, es un mamífero» es verdadera; su recíproco («si es mamífero, es perro») es falso. Su contrarrecíproco
(«si no es mamífero, no es perro») es tan verdadero como la original.

\subsubsection{Las dos inferencias válidas y las dos trampas}
\begin{itemize}
  \item \clave{Modus ponens}: $p \to q$ y $p$; luego $q$. ✓
  \item \clave{Modus tollens}: $p \to q$ y $\neg q$; luego $\neg p$. ✓
  \item Afirmar el consecuente: $p \to q$ y $q$; luego $p$. ✗ (el partido se pudo cancelar por otra razón)
  \item Negar el antecedente: $p \to q$ y $\neg p$; luego $\neg q$. ✗
\end{itemize}

\subsubsection{Otras formas de decir «si…, entonces»}
«$p$ solo si $q$» es $p \to q$. «$q$ si $p$» es $p \to q$. «$p$ es suficiente para $q$» es $p \to q$; «$q$ es necesario
para $p$» también. Y $p \to q \equiv \neg p \vee q$: «si estudias, apruebas» equivale a «no estudias o apruebas».

\figura{cartas}{Las cuatro cartas: solo dos pueden romper la regla.}

\begin{regla}
\begin{itemize}
  \item $p \to q \equiv \neg q \to \neg p \equiv \neg p \vee q$.
  \item Válidas: modus ponens y modus tollens. Inválidas: afirmar el consecuente y negar el antecedente.
  \item «$p$ solo si $q$»: $p \to q$. «Suficiente» va a la izquierda de la flecha; «necesario», a la derecha.
\end{itemize}
\end{regla}

\subsection{Ejemplos resueltos}
\begin{ejemplo}[Las cuatro cartas]
Cada carta tiene una letra de un lado y un número del otro. Regla: «si una carta tiene vocal, al otro lado tiene número
par». Se ven: \fbox{E}\ \fbox{K}\ \fbox{4}\ \fbox{7}. ¿Qué cartas hay que voltear para comprobar la regla?
\solucion La E (modus ponens: debe tener par) y el 7 (modus tollens: si tuviera vocal, rompería la regla). La K no
importa y el 4 tampoco (afirmar el consecuente no prueba nada). Se voltean E y 7.
\end{ejemplo}

\begin{ejemplo}[Solo si]
«Te dan la beca solo si tu promedio es mayor que 4.5». Carlos tiene promedio 4.7. ¿Le dan la beca?
\solucion La frase es «beca $\to$ promedio $> 4.5$». Saber el consecuente no permite concluir el antecedente: no se puede
asegurar. Lo que sí se sabe: con promedio 4.3 no habría beca.
\end{ejemplo}

\begin{trampa}
Afirmar el consecuente es el error más común en las preguntas de lógica del examen. Si la conclusión «vuelve» por la
flecha al revés, desconfía.
\end{trampa}

\taller
\piso{Entrena}{condicionales}
\begin{ejercicios}
\item Escribe el contrarrecíproco de «si estudio, apruebo». \sol{«Si no apruebo, no estudié».}
\item «Si es cuadrado, es rectángulo». Una figura es rectángulo. ¿Es cuadrado? \sol{No se sabe (afirmar el consecuente).}
\item Traduce: «Aprobar es necesario para graduarse». \sol{«Si me gradúo, aprobé»: graduarse $\to$ aprobar.}
\end{ejercicios}
\piso{Como en el examen}{opción múltiple}
\begin{ejercicios}
\item ¿Cuál equivale a «si un número es múltiplo de 6, es par»?
\opcionesl{Si un número es par, es múltiplo de 6.}{Si un número no es par, no es múltiplo de 6.}{Si un número no es múltiplo de 6, no es par.}{Un número es par y múltiplo de 6.}
\sol{B. Contrarrecíproco.}
\item «Si hay huelga, no hay clase. Si no hay clase, voy a la biblioteca.» Hoy no fui a la biblioteca. Entonces
\opcionesd{hubo huelga}{hubo clase y no hubo huelga}{no hubo clase}{hubo huelga y clase}
\sol{B. Por tollens: hubo clase; y de nuevo por tollens: no hubo huelga.}
\end{ejercicios}

\begin{retoresuelto}
Es modus tollens: lluvia $\to$ cancelación; no hubo cancelación; luego no llovió. \textbf{Opción B.}
\end{retoresuelto}
\bitacora

% ------------------------------------------------------------------
\leccion{pu-lo-4}{Cuantificadores y silogismos}{Juego de los diagramas que encierran.}

\begin{reto}
«Todos los ingenieros saben cálculo. Algunos músicos son ingenieros.» De estas premisas se concluye que
\opcionesl{Algunos músicos saben cálculo.}{Todos los músicos saben cálculo.}{Algunos que saben cálculo no son músicos.}{Ningún músico sabe cálculo.}
\end{reto}

\subsection{La idea: dibujar conjuntos}
Las frases con «todos», «algunos» y «ningún» hablan de conjuntos. Dibujarlas con círculos (diagramas de Euler) vuelve
evidente qué se sigue y qué no:
\begin{itemize}
  \item «Todos los $A$ son $B$»: el círculo $A$ dentro del círculo $B$.
  \item «Ningún $A$ es $B$»: círculos separados.
  \item «Algunos $A$ son $B$»: círculos que se cruzan, con una marca $\times$ en la parte común (hay al menos uno ahí).
\end{itemize}

\figura{euler}{Las premisas dibujadas: basta un músico dentro del círculo de ingenieros.}

\subsubsection{Cómo se prueba que una conclusión \emph{no} se sigue}
Basta dibujar un caso que cumpla las premisas y donde la conclusión falle. Si no puedes dibujar un contraejemplo de
ninguna manera, la conclusión es válida.

\begin{regla}
\begin{itemize}
  \item Una conclusión es válida si es verdadera en \emph{todos} los dibujos que respetan las premisas.
  \item «Todos los $A$ son $B$» no dice que existan $A$, ni que todos los $B$ sean $A$.
  \item De dos premisas con «algunos» no se concluye nada seguro.
\end{itemize}
\end{regla}

\subsection{Ejemplos resueltos}
\begin{ejemplo}[Un silogismo inválido]
«Todos los gatos son mamíferos. Algunos mamíferos nadan. Luego algunos gatos nadan.» ¿Es válido?
\solucion No: se puede dibujar el círculo de gatos dentro de mamíferos y los que nadan en otra zona de mamíferos, sin
tocar a los gatos.
\end{ejemplo}

\begin{ejemplo}[Ninguno]
«Ningún reptil tiene pelo. Todos los perros tienen pelo.» ¿Qué se concluye?
\solucion El círculo de perros está dentro de «tienen pelo», que no toca a los reptiles: ningún perro es reptil.
\end{ejemplo}

\begin{trampa}
Concluir lo que es cierto en la vida real pero no se sigue de las premisas. El examen evalúa la \emph{forma} del
argumento, no si la conclusión es verdad en el mundo.
\end{trampa}

\taller
\piso{Entrena}{dibujar}
\begin{ejercicios}
\item «Todos los $A$ son $B$. Todos los $B$ son $C$.» ¿Qué se concluye de $A$ y $C$? \sol{Todos los $A$ son $C$.}
\item «Algunos $A$ son $B$. Algunos $B$ son $C$.» ¿Se concluye que algunos $A$ son $C$? \sol{No: los $B$ que son $A$ pueden ser otros que los que son $C$.}
\item «Ningún $A$ es $B$. Algunos $C$ son $A$.» ¿Qué se concluye? \sol{Algunos $C$ no son $B$ (los que son $A$).}
\end{ejercicios}
\piso{Como en el examen}{opción múltiple}
\begin{ejercicios}
\item «Todos los deportistas del club madrugan. Lucía no madruga.» Entonces
\opcionesd{Lucía es del club.}{Lucía no es deportista del club.}{Lucía es deportista.}{No se puede concluir nada.}
\sol{B. Si fuera deportista del club, madrugaría.}
\item «Ningún político es honesto» y «algunos abogados son políticos». Se concluye:
\opcionesl{Algunos abogados no son honestos.}{Ningún abogado es honesto.}{Algunos honestos son abogados.}{Todos los abogados son políticos.}
\sol{A. Los abogados que son políticos no son honestos.}
\end{ejercicios}

\begin{retoresuelto}
El círculo de ingenieros está dentro del de «saben cálculo», y hay al menos un músico en el círculo de ingenieros. Ese
músico sabe cálculo. \textbf{Opción A.} La C puede ser cierta pero no se sigue: quizás todos los que saben cálculo son
músicos.
\end{retoresuelto}
\bitacora

% ------------------------------------------------------------------
\leccion{pu-lo-5}{Verdad y mentira}{Juego de los guardianes que dicen la verdad o mienten.}

\begin{reto}
Tres personas, A, B y C, son veraces (siempre dicen la verdad) o mentirosas (siempre mienten). A dice: «B miente». B dice:
«C miente». C dice: «A y B mienten». ¿Quiénes dicen la verdad?
\opciones{solo A}{solo B}{A y C}{ninguno}
\end{reto}

\subsection{La idea: suponer y buscar la contradicción}
En estos acertijos se \clave{supone} algo («A dice la verdad») y se sigue la cadena de consecuencias. Si se llega a una
contradicción, la suposición era falsa. Si no, se tiene una solución (que conviene verificar con todas las frases).

\subsubsection{Afirmaciones que se contradicen}
Si dos personas dicen cosas \emph{contradictorias} (una es la negación de la otra), exactamente una dice la verdad. Es el
mejor punto de partida: ya se sabe algo sin suponer nada.

\subsubsection{Contar mentirosos}
Cuando el enunciado dice «solo uno miente», se prueba cada persona como la mentirosa y se ve cuál deja todo coherente.

\begin{regla}
\begin{enumerate}
  \item Busca dos frases contradictorias: una es verdad y otra mentira.
  \item Si no las hay, supón el valor de la persona con la frase más «fuerte» y encadena.
  \item Verifica la solución con todas las frases.
\end{enumerate}
\end{regla}

\subsection{Ejemplos resueltos}
\begin{ejemplo}[Solo uno miente]
Se rompió un vidrio. Ana: «Fue Beto». Beto: «Fue Ciro». Ciro: «Beto miente». Dani: «Yo no fui». Solo uno miente. ¿Quién
rompió el vidrio?
\solucion Beto y Ciro se contradicen: uno de ellos es el mentiroso, así que Ana y Dani dicen la verdad: fue Beto (y Dani
no fue). Entonces Beto miente al decir «fue Ciro», y Ciro dice la verdad. Fue Beto.
\end{ejemplo}

\begin{ejemplo}[El guardián]
Un guardián dice: «Yo soy mentiroso». ¿Qué es?
\solucion Si fuera veraz, diría una mentira; si fuera mentiroso, diría una verdad. Ninguna opción funciona: la frase es una
paradoja y nadie de ese mundo puede decirla. El examen usa esto para descartar opciones.
\end{ejemplo}

\begin{trampa}
Olvidar verificar la solución con \emph{todas} las frases. A veces una suposición no lleva a contradicción en dos frases
pero choca con la tercera.
\end{trampa}

\taller
\piso{Entrena}{suponer}
\begin{ejercicios}
\item X dice: «Y y yo somos del mismo tipo». Y dice: «X es mentiroso». ¿Qué es cada uno? \sol{Si X veraz, Y veraz, pero Y diría una mentira: contradicción. X mentiroso: entonces son de distinto tipo, Y veraz, y Y dice la verdad ✓.}
\item Cuatro sospechosos; solo uno dice la verdad. P: «Fue Q». Q: «Fue S». R: «Yo no fui». S: «Q miente». ¿Quién fue? \sol{Q y S se contradicen: uno dice la verdad. Entonces P y R mienten: no fue Q y R sí fue. Fue R (y Q miente, S dice la verdad).}
\end{ejercicios}
\piso{Como en el examen}{opción múltiple}
\begin{ejercicios}
\item Tres hermanos, uno de los cuales se comió el pastel. Juan: «Fui yo». Kike: «No fue Juan». Leo: «No fui yo». Si exactamente uno dice la verdad, ¿quién se comió el pastel?
\opciones{Juan}{Kike}{Leo}{no se puede saber}
\sol{C. Juan y Kike se contradicen: uno dice la verdad. Así que Leo miente: fue Leo. Verificación: Juan miente y Kike dice la verdad.}
\end{ejercicios}

\begin{retoresuelto}
Si A dice la verdad: B miente, entonces C dice la verdad, y C afirma que A miente: contradicción. Luego A miente: B dice la
verdad, C miente, y la frase de C («A y B mienten») es falsa porque B es veraz ✓. Solo B dice la verdad.
\textbf{Opción B.}
\end{retoresuelto}
\bitacora

% ------------------------------------------------------------------
\evaluacion
\begin{ejercicios}
\item Si $p$ es F y $q$ es V, el valor de $(p \to q) \wedge (q \to p)$ es
\opciones{V}{F}{depende de $r$}{no se puede evaluar}
\sol{B. V $\wedge$ F. (2.1)}
\item Si $(p \vee q) \to r$ es falsa y $p$ es F, entonces
\opciones{$q$ es V y $r$ es F}{$q$ es F y $r$ es V}{$q$ y $r$ son F}{$q$ y $r$ son V}
\sol{A. Antecedente V (luego $q$ V) y consecuente F. (2.1)}
\item La negación de «Todos los días hace frío o llueve» es
\opcionesl{Algún día no hace frío y no llueve.}{Ningún día hace frío ni llueve.}{Algún día no hace frío o no llueve.}{Todos los días hace calor y no llueve.}
\sol{A. (2.2)}
\item $\neg(p \to \neg q)$ equivale a
\opciones{$p \wedge q$}{$\neg p \wedge q$}{$p \vee q$}{$p \wedge \neg q$}
\sol{A. $\neg(p \to r) = p \wedge \neg r$ con $r = \neg q$. (2.2)}
\item «Si como dulces, me duele la muela.» ¿Cuál es equivalente?
\opcionesl{Si me duele la muela, comí dulces.}{Si no como dulces, no me duele la muela.}{Si no me duele la muela, no comí dulces.}{Como dulces y me duele la muela.}
\sol{C. (2.3)}
\item «Solo los socios pueden entrar.» Marta entró. Entonces
\opcionesd{Marta es socia.}{Marta no es socia.}{No se sabe.}{Todos los socios entraron.}
\sol{A. Entrar $\to$ socio. (2.3)}
\item «Algunos artistas son tímidos. Todos los tímidos son callados.» Se concluye que
\opcionesl{Algunos artistas son callados.}{Todos los artistas son callados.}{Algunos callados no son artistas.}{Ningún artista es callado.}
\sol{A. (2.4)}
\item Dos guardianes: uno siempre dice la verdad y otro siempre miente. El primero dice: «Los dos decimos la verdad». ¿Quién es el veraz?
\opciones{el primero}{el segundo}{ninguno}{no se sabe}
\sol{B. Si el primero fuera veraz, la frase sería falsa. (2.5)}
\end{ejercicios}
